\chapter{Examples and complements}
\label{XI}
% original p. 285

\section{Projective spaces, unirational varieties}
\label{XI.1}

\begin{proposition}
\label{XI.1.1}
Let $k$ be an algebraically closed field, $X=\PP_k^r$ the projective
space of dimension $r$ over~$k$. Then $X$ is \emph{simply
connected,} \ie $\pi_1(X)=0$.
\end{proposition}

For $r=0$, this is trivial. If $r=1$, one must show that if $X'$ is
a nonempty connected \'etale covering of~$X=\PP_k^1$, then
$X'\isomto X$. The genus formula gives us here, if~$g$ and $g'$ are
the genera of~$X$ and~$X'$:
$$
1-g' = d(1-g)\quoi,
$$
where $d$ is the degree of~$X'$ over~$X$. Since $g=0$, we will therefore have
$1-g'=d$, which requires $d=1$ since $g'\ge 0$, which proves
$X'\isomto X$. When $r\ge 2$, one proceeds by induction
on~$r$, assuming that $\PP^{r'}$ is simply connected for
$r'<r$. Applying this to a hyperplane of~$\PP^r$ and using
(X~\Ref{X.2.10}), it indeed follows that $\PP^r$ is simply
connected. Another proof: we will have
$\pi_1(\PP^1\times\cdots\times \PP^1) = \pi_1(\PP^1)\times\cdots
\times \pi_1(\PP^1)$ by virtue of (X~\Ref{X.1.7}), hence $(\PP^1)^r$
is simply connected since $\PP^1$ is, hence $\PP^r$ is
simply connected by virtue of the birational invariance of the fundamental
group (X~\Ref{X.3.4}). This proof shows more
generally:

\begin{corollary}
\label{XI.1.2}
Let $X$ be a proper and normal scheme over an
algebraically closed field~$k$; if $X$ is a rational
variety, \ie integral and its function field is a purely
transcendental extension of~$k$, then $X$ is simply connected.
\end{corollary}

This result applies in particular to Grassmannian
varieties and more
% original p. 286
generally to the varieties~$G/H$, where $G$ is a connected linear
group over~$k$ and $H$ an algebraic subgroup containing a Borel
subgroup of~$G$.

Recall that one calls \emph{unirational variety over~$k$}
a proper and integral scheme over~$k$ whose function
field $K$ is contained in a purely transcendental extension~$K'$
of~$k$, finite over~$K$ (\ie having the same transcendence degree
over~$k$ as~$K$), \ie if there exists a dominant rational
map $f\colon\PP_k^r\to X$, with $r=\dim X$. If $X$ is normal,
one therefore sees by the reflections preceding X~\Ref{X.3.4}
that for every connected \'etale covering $X'$ of~$X$, with
field~$L/K$, the algebra $L\otimes_K K'$ over~$K'$ is
unramified over the model~$\PP^r$, hence completely
decomposed by virtue of~\Ref{X.1.1}, which shows that $L$ is
$K$-isomorphic to a subextension of~$K'/K$. This therefore proves,
taking into account V~\Ref{V.8.2}:

\begin{corollary}
\label{XI.1.3}
The fundamental group of a normal unirational variety over
an algebraically closed field is finite.
\end{corollary}

(N.B. One will note that in the definition of ``unirational'', one
did not need $K'/K$ to be finite).

\begin{remarks}
\label{XI.1.4}
Of course, the results of this no.\ are well
known. On the other hand,
J-P\ptbl
Serre has shown~\cite{XI.10} that when $X$ is
a smooth unirational projective variety over an
algebraically closed field of \emph{characteristic zero,} $X$ is
simply connected. His proof is transcendental in that
it uses the Hodge symmetry theorem.

\lcrochetbf Added in 2003 (MR) Let us restrict ourselves to the case of an algebraically closed field $k$ of
characteristic $p\geq 0$. In characteristic $p>0$, the definition of
unirational $k$-variety given above is that of weakly unirational
$k$-variety, as opposed to that of strongly unirational
$k$-variety, where one assumes moreover that $K'$ is a separable
extension of~$K$.

In the case of dimension $2$, there exist projective and
smooth, weakly unirational surfaces which have a nontrivial fundamental group (finite
by 1.3) and which are therefore not rational (T\ptbl Shioda, On unirationality of supersingular surfaces, Math.~Ann. \textbf{225}
(1977), p\ptbl155--159.). On the other hand, a strongly
unirational surface is rational: this follows from Castelnuovo's
rationality criterion, extended to characteristic $p>0$ by O\ptbl Zariski (\cf J-P\ptbl Serre, S\'em. Bourbaki \no146, 1957 and {\OE}uvres--Collected
papers, vol\ptbl 1, p\ptbl 467).

Over the field $\CC$ of complex numbers, one knows examples of smooth projective
varieties of \hbox{dimension $\geq 3$} which are unirational and not
rational (\cf P\ptbl Deligne, Vari\'et\'es unirationnelles non
rationnelles [d'apr\`es M\ptbl Artin et D\ptbl Mumford],  S\'em.\ Bourbaki
\no402, 1971-72, Lecture Notes vol\ptbl 317). This is the case in particular for a
smooth cubic hypersurface in the projective space $\PP^4$.
Some of these examples extend to characteristic $>0$ to give
strongly unirational non-rational varieties (\cf J.P\ptbl Murre, Reduction of the proof of the non-rationality of a non
singular cubic threefold to a result of Mumford, Compositio Math. \textbf{27}
(1973), p\ptbl63--82).

Let $V$ be a normal, integral projective $k$-variety. One says that $V$ is
rationally connected if there exist an integral $k$-scheme of finite type
$T$ and a $k$-morphism $F : T\times\PP^1\to V$ such that the morphism:
\begin{equation}\tag*{(1)}
\begin{split}
T \times\PP^1\times\PP^1&\to V\times V\\
(t, u, u')&\mto(F(t,u), F(t,u'))
\end{split}
\end{equation}
is dominant. In particular, through two sufficiently general rational
points of~$V$ there passes a rational curve. The terminology is
justified by the fact that if $V$ is rationally connected, one can
join two rational points by a finite chain of rational
curves. If $k$ is of characteristic $p>0$, one strengthens the
preceding definition by requiring that the map (1) be
generically smooth. One then obtains the notion of separably
rationally connected variety. Thus unirational varieties
are rationally connected and in characteristic
$p>0$, strongly unirational varieties are separably
rationally connected. J\ptbl Koll\'ar has shown that separably
rationally connected varieties have a trivial algebraic
fundamental group (\cf O\ptbl Debarre, Vari\'et\'es rationnellement
connexes [d'apr\`es T\ptbl Graber, J\ptbl Harris, J\ptbl Starr et A.J\ptbl de~Jong],
S\'em. Bourbaki \no906, 2001-2002).\rcrochetbf
\end{remarks}

\section{Abelian varieties}
\label{XI.2}
Let $k$ be an algebraically closed field, $A$ an abelian
variety over~$k$, \ie a group scheme over~$k$, proper
over~$k$, smooth over~$k$, and connected, and finally $G$ a
commutative group scheme of finite type over~$k$. Let us denote by $\Ext(A,G)$ the
group of classes of commutative extensions of~$A$ by~$G$, by
$\H^1(A,G)$ the group of classes of principal bundles on~$A$ with
group~$G$ (compare \No \Ref{XI.4} below), and let us consider
the canonical homomorphism
$$
\Ext(A,G)\to \H^1(A,G)\quoi.
$$
An
% original p. 287
argument of Serre \cite[Chap\ptbl VII, Th\ptbl 5]{XI.5} shows that it is an
injective homomorphism, whose image is the set of
``primitive elements'' of~$\H^1(A,G)$, \ie of the
elements $\xi$ for which one has:
$$
\pi^*(\xi) = \pr_1^*(\xi) + \pr_2^*(\xi)\quoi,
$$
where $\pr_i$ are the two projections of~$A\times A$
onto~$A$, and $\pi\colon A\times A\to A$ the composition law of~$A$
(N.B. Serre states his theorem only for $G$
linear and connected, and of course smooth over~$k$, but by
simplifying the first part of his argument, one sees that
these restrictions are unnecessary: it suffices to note that every morphism
from~$A$ to a group scheme $E$ of finite type over~$k$ which
takes the unit to the unit is a group homomorphism, and
to apply this to the sections over~$A$ of an extension $E$ of
$A$ by~$G$).

We are going to apply this result to the case where $G$ is a finite
group separable over~$k$, \ie an ordinary finite group, assumed
commutative. Using then $\pi_1(A\times A)\isomto\allowbreak \pi_1(A)\times
\pi_1(A)$ (X~\Ref{X.1.7}) and interpreting $\H^1(X,G)$ as
$\Hom(\pi_1(X),G)$ for every algebraic scheme~$X$, in
particular for $X=A$ or $X=A\times A$, one sees that every class of
$\H^1(A,G)$ is primitive, hence one has an isomorphism
$$
\Ext(A,G)\isomto \H^1(A,G)\quoi,
$$
in other words \emph{every principal covering of~$A$ with
commutative structure group~$G$, pointed over the origin
of~$A$, is endowed in a unique way with a structure of
algebraic group admitting the marked point as origin, and such that
$A'\to A$ is a homomorphism of algebraic groups.} In
particular, if $A'$ is connected, it is also an
abelian variety, isogenous to~$A$.

On the other hand, since the functor $X\mto \pi_1(X)$ from pointed
algebraic schemes $X$ to groups commutes with products
(IX~\Ref{IX.1.7}), it transforms a group in the first
category into a group in the category of groups, \ie into
a \emph{commutative} group. Hence \emph{if $A$ is an abelian
variety, $\pi_1(A)$
% original p. 288
is a commutative group.} Hence, in order to know $\pi_1(A)$, it
suffices to know the functor $G\mto
\H^1(A,G)=\Hom(\pi_1(A),G)$ for $G$ varying over the
\emph{commutative} finite groups. Finally, let us recall that for every integer $n>0$,
the homomorphism of multiplication by $n$ in~$A$:
$$
A\lto{n} A
$$
is surjective, hence with finite kernel, \ie it is an isogeny, and
that it follows that every isogeny $A'\to A$ is a quotient
of an isogeny of the preceding type. From this, and from
standard arguments (\cf for example~\cite{XI.6}), one derives:

\begin{theorem}[Serre-Lang]
\label{XI.2.1}
Let $A$ be an abelian variety over an
algebraically closed field~$k$, and for every integer $n>0$ let us consider
the ordinary finite group $K_n$ underlying the kernel ${}_nA$ of the
multiplication by $n$ in~$A$; finally, let us set, for every prime
number~$\ell$:
$$
T_\ell(A)=\varprojlim_r K_{\ell^r}
$$
and
$$
T_{.}(A)=\underset{\ell}{\prod}T_\ell(A)=\varprojlim_n K_n
$$
(where, for $m$ a multiple of~$n$, $m=ns$, one maps $K_m$ into $K_n$
by multiplication by~$s$). Then the group $\pi_1(A)$ is
canonically isomorphic to $T_{.}(A)$, hence for every prime
number $\ell$, the $\ell$-primary component of~$\pi_1(A)$ is
canonically isomorphic to $T_\ell(A)$.
\end{theorem}

One will note that these isomorphisms are functorial for variable $A$. The
module $T_\ell(A)$ is called the \emph{$\ell$-adic Tate
module} of the abelian variety~$A$. It is an
additive functor in~$A$; in particular it gives rise to a
representation of the ring $\Hom(A,A)$ of endomorphisms of~$A$
in $T_\ell(A)$, called the \emph{$\ell$-adic representation
of Weil}, and which plays an important role in the theory of
abelian varieties (\cf for example \hbox{\cite[Chap~VII]{XI.4}}).
Theorem~\Ref{XI.2.1} gives an interpretation of it in
terms of the natural representation in the \emph{$\ell$-adic homology
group} of~$A$, $\H_1(A,\ZZ_\ell)=\pi_1(A)_\ell$,
which
% original p. 289
is evidently more satisfying a priori, notably from the point of view
of the Lefschetz formula, \cite[Chap~V]{XI.4}. Let us recall here
Weil's results on the structure of~$T_\ell(A)$:
\begin{enumerate}
\item[a)] If $n$ is prime to $\car(k)$, then $K_n$ is a free
module of rank equal to $2\dim A$ over~$\ZZ/n\ZZ$, hence if $\ell$
is a prime number $\neq \car(k)$, $T_\ell(A)$ is a free module
of rank equal to $2\dim A$ over the ring $\ZZ_\ell$ of
$\ell$-adic integers;
\item[b)] If $n$ is a power of~$\car(k)=p$, then $K_n$ is a
free module of rank $\nu\leq\dim A$ over~$\ZZ/n\ZZ$, $\nu$
independent of~$n$, hence $T_p(A)$ is a free module of rank
$\nu\leq \dim A$ over the ring $\ZZ_p$ of $p$-adic integers.
\end{enumerate}

This shows that in the theory of the fundamental group
developed here, the fundamental group of a variable abelian
variety does not vary in a regular way
with the parameter, its $p$-primary component being liable to decrease
abruptly for values of the parameter $t$ corresponding to
a residue characteristic~$p$; the best-known case of
this phenomenon is that of elliptic curves. One will note,
however, that whatever $n$ may be (prime or not to the
characteristic), the true kernel ${}_nA$ in $A$ for the
multiplication by $n$ is a finite group scheme over~$k$ of
degree $n^{2g}$, where $g=\dim A$, which will be non-separable over~$k$ if $n$ is a multiple of~$p=\car(k)$. Moreover, when $A$
varies in a family of abelian varieties, \ie if one
has an abelian scheme $A$ over a base scheme $S$, one
shows more generally that ${}_n A$ is a finite and flat group
scheme over~$S$, of degree $n^{2g}$ over~$S$;
that is to say, provided one takes into account the infinitesimal
parts of the kernels~${}_n A$, they behave in a regular way
whatever $n$ may be. This suggests that the
``true'' fundamental group of an abelian variety $A$
is the algebraic pro-group (formal projective limit of
finite groups over~$k$) $\varprojlim_n{{}_nA}$, where
by ``true fundamental group'' of an algebraic scheme~$X$, one
must understand: the pro-group which classifies the principal
coverings of~$X$ with structure group an arbitrary finite group $G$
over~$k$ (not necessarily separable over~$k$). In this
way, for example, one recovers, through the
representations of~$\Hom(A,A)$ in the $p$-primary component of the
true fundamental group of~$A$, the characteristic polynomial of
Weil defined by the latter with the help of the $\ell\neq p$, in a more natural way than Serre's construction~\cite{XI.8}.

\section{Projecting cones, Zariski's example}
\label{XI.3}
% original p. 290

As before, let $k$ be an algebraically closed field for simplicity, and
let $V$ be a connected projective $k$-scheme, a closed
subscheme of~$\PP^r_k$, which one may, if one wishes, assume
nonsingular. Let $Y=\hat C$ be the projective projecting cone of~$V$,
$y_0$ its vertex, $X=\hat C_V$ the usual projective closure of the
vector bundle $C_V=\VV(\cal{O}_V(1))$ associated with
$\cal{O}_V(1)$, and finally
$$
f\colon X\to Y
$$
the canonical morphism, contracting the zero section $X_0$ of~$C_V$
over~$X$ to a point (EGA~II~8.6.4). Since $X$ is a locally
trivial bundle over~$V$, with fibers $\PP^1$, hence with
simply connected fibers, the morphism $p\colon X\to V$ induces, by virtue
of~(X~\Ref{X.1.4}),
an isomorphism:
$$
\pi_1(X)\isomto\pi_1(V)\quoi.
$$
Since $p$ induces an isomorphism $X_0\to V$, one concludes that \emph{an
\'etale covering of~$X$ is completely decomposed
if and only if its restriction to $X_0$ is}. Now for every
\'etale covering $Y'$ of~$Y$, the inverse image $X'=X\times_Y
Y'$ is an \'etale covering of~$X$ completely
decomposed over the fiber~$X_0$, hence trivial. Since
the homomorphism $\pi_1(X)\to\pi_1(Y)$ is surjective (IX~\Ref{IX.3.4}),
one concludes that
$$
\pi_1(Y)=(e)
$$
in other words \emph{every projective projecting cone is
simply connected.} (N.B. in characteristic~$0$, the
same result will be valid if one takes for $Y$ the affine projecting
cone).

Suppose now $V$ regular, \ie smooth over~$k$; then
$X$ is regular, and for a suitable projective embedding
of~$V$, one then finds a \emph{normal} projecting cone $Y$. If
$V$ is not simply connected, hence $X$ not simply connected,
let $X'$ be a nontrivial connected \'etale covering
of~$X$. Since the fibers of~$X$ at the points $y\in Y$ distinct from
$y_0$ are reduced to a point, one sees that the restriction of
$X'$ to its fibers (in particular to the
generic
% original p. 291
fiber) is trivial; however $X$ does not come by
inverse image from an \'etale covering of~$Y$, since $Y$ is
simply connected and $X'$ would be completely
decomposed. This shows that (X~\Ref{X.1.3} and~\Ref{X.1.4})
become false if one replaces the hypothesis that $f$ is
separable by the weaker one that its fibers are separable
(or even smooth) algebraic schemes over
the~$\kres(s)$. One will note likewise that the fundamental groups of the
geometric fibers $\overline X_y$ of the $y\neq y_0$ are
evidently reduced to $(e)$ since these fibers are
reduced to a point, whereas $\pi_1(X_0)\neq e$; hence the
semicontinuity theorem (X~\Ref{X.2.4}) likewise
fails for~$f$.

Let us finally indicate the example, pointed out by Zariski, which shows
that these same theorems fail when one replaces in them
the hypothesis that $f$ is separable by the hypothesis that $f$ is
flat. Let $f\colon X\to Y$ be a morphism from a nonsingular projective
surface to the rational line $Y=\PP^1$, such
that
$K=k(X)$
is a ``regular'' extension, \ie primary and
separable, of~$k(f)$ (\ie the geometric generic
fiber is connected and separable), and such that the
divisor $(f)=X_0-X_\infty$ is an
$n$th multiple
of a divisor (where $n$ is an integer prime to the
characteristic). It is possible to construct such examples
in every characteristic. Let $X'$ be the normalization of~$X$ in
$K(f^{1/n})$, where $K=k(X)$ is the function field of~$X$. It
follows from the hypothesis on~$(f)$ that $X'$ is \'etale
over~$X$. Let $Y'$ be the normalization of~$Y$ in $k(t)(t^{1/n})$; it
is ramified over~$Y$ at the points $t=0$ and $t=\infty$ exactly,
and the restriction $X'|f^{-1}(U)$ is isomorphic to the inverse image
of~$Y'|U$. In particular, the restriction of~$X'$ to the
\emph{geometric} generic fiber of~$X$ over~$Y$
decomposes completely. However, $X'$ is not isomorphic
to the inverse image of an \'etale covering of~$Y$, for one
sees at once that the latter would necessarily be~$Y'$, which
is absurd since $Y'$ is ramified over~$Y$\kern2pt\footnote{One may
remark, from the point of view of the ``\'etale topology'' (SGA~4~VII),
that in this example $\R^1(f_{\et})_*(\ZZ/n\ZZ)$ is ``non-separated''
over~$S$.}.

Here is (after Serre) a simple way of realizing the
conditions of this example, drawing inspiration from \cite[\No 20]{XI.5}: one takes
for $n$ a prime number $\geq 5$, distinct from the
characteristic, and one lets $G=\ZZ/n\ZZ$ act on
$\underline{k}^4$%
{{\renewcommand{\thefootnote}{*}\addtocounter{footnote}{-1}\kern1pt\footnote{\lcrochetbf Added in 2003: $\underline{k}^4$ denotes the affine space of dimension $4$ over $k$.\rcrochetbf}}}
by multiplying the coordinates by four distinct
characters of~$G$ (which is possible since $n\geq
5$). Then $G$ acts on the projective space $\PP^3_k$, and the
only points fixed under $G$ are the four points corresponding to the
coordinate axes. The surface $X'$ with equation
$x^n+y^n+z^n+t^n=0$
is smooth over
% original p. 292
$k$ (Jacobian criterion), and contains none of the fixed points;
hence $G$, being of prime order, acts on~$X$ ``without fixed
points'', \ie $X$ is a principal covering of~$X=X'/G$ with group
$G$. Let $g=x/y$ in $k(X')=K'$; it is a Kummer
generator of~$K'$ over~$K=k(X)$ if the chosen characters
were~$\chi^i$, $i=0,1,2,3$, with $\chi$ a primitive
character; let $f$ be its $n$th power, which is an
element of~$K$. One sees at once that $K'$ is a
regular extension of~$k(g)$, which follows from the fact that
the plane curve with homogeneous equation in $U,T,Z$:
$T^n+Z^n+(1+g^n)U^n=0$ is smooth over~$k(g)$ (Jacobian criterion), and
that one knows that every plane curve is connected. On the other hand one has
$k(f)=K\cap k(g)$, since the right-hand side is an extension of
$k(f)$ contained in the extension $k(g)$ of prime degree, and
distinct from~$k(g)$ (since $g\not\in K$). This implies that $K$ is
a regular extension of~$k(f)$. Finally the divisor of~$f$
on~$X$ is an $n$th multiple of a divisor, since its
inverse image on~$X'$ is the divisor of~$g^n$, hence an $n$th
multiple, and one can descend back because $X'$ is \'etale over~$X$. One would be done if the rational map $f\colon X\to \PP^1$
were a morphism, that is to say if the divisors of
zeros and of poles of~$f$ did not meet. In fact,
one easily verifies (again by looking on~$X'$) that the
two divisors in question are the products by $n$ of two
curves smooth over~$k$, intersecting transversally at a point~$a$. Replacing now $X$ by the scheme obtained by blowing up~$a$, the preceding conditions ($\divisor(f)$ divisible by
$n$, and $k(X_1)=k(X)$ regular extension of~$k(X)$) remain
satisfied, but moreover $f$ is a \emph{morphism} $X_1\to
\PP^1$, hence we are under the desired conditions.
